\chapter{Обзор литературы}\label{ch:lit_review}

В данной главе представлен обзор основных направлений исследований, связанных с проблемами безопасности больших языковых моделей в мультимодальном сценарии. Рассматриваются ключевые технологии, лежащие в основе современных систем обработки естественного языка: текстовые эмбеддинги, методы переранжирования, подходы генерации с извлечением информации (RAG), а также системы обеспечения безопасности и уязвимости мультимодальных моделей.

%% =================================================================
\section{Эмбеддинги текста}\label{sec:embeddings}
%% =================================================================

Текстовые эмбеддинги --- векторные представления текстовых фрагментов в непрерывном пространстве --- являются фундаментальным компонентом современных систем информационного поиска, кластеризации и классификации. Развитие данного направления прошло путь от классических методов, основанных на архитектуре BERT, до современных подходов, использующих большие языковые модели в качестве основы для получения универсальных представлений.

\subsection{Классические методы}\label{subsec:classic_embeddings}

Одной из первых работ, предложивших эффективный метод получения предложенческих эмбеддингов, стала Sentence-BERT (SBERT)~\cite{1908.10084}. Авторы использовали сиамскую и триплетную архитектуры поверх предобученной модели BERT для получения семантически осмысленных представлений предложений, допускающих сравнение с помощью косинусного сходства. Данный подход сократил время поиска наиболее похожей пары среди 10\,000 предложений с 65 часов до примерно 5 секунд при сохранении точности уровня BERT.

SimCSE~\cite{2104.08821} предложил простой, но эффективный метод контрастивного обучения предложенческих эмбеддингов. В ненаблюдаемом варианте модель использует dropout в качестве минимального возмущения данных: одно и то же предложение пропускается через энкодер дважды с различными масками dropout, образуя позитивные пары. Данный подход продемонстрировал значительное улучшение качества на задачах семантического текстового сходства.

DiffCSE~\cite{2204.10298} развил идеи SimCSE, предложив контрастивное обучение, основанное на различиях: модель обучается быть чувствительной к разнице между исходным предложением и его отредактированной версией, что позволяет лучше улавливать семантические нюансы.

\subsection{Модели на основе больших языковых моделей}\label{subsec:llm_embeddings}

С ростом масштаба языковых моделей исследователи обнаружили, что декодерные архитектуры, изначально предназначенные для генерации текста, могут быть успешно адаптированы для построения высококачественных эмбеддингов.

E5~\cite{2212.03533} представил семейство моделей текстовых эмбеддингов, обучаемых контрастивным методом на большом корпусе слабо размеченных текстовых пар (CCPairs). E5 стала первой моделью, превзошедшей сильную базовую линию BM25 на бенчмарке BEIR без использования размеченных данных. Дальнейшее развитие --- E5-Mistral~\cite{2401.00368} --- предложило получать высококачественные текстовые эмбеддинги с использованием исключительно синтетических данных, сгенерированных большими языковыми моделями для сотен тысяч задач на 93 языках, с последующей дообучкой декодерных моделей с открытым кодом.

GTE~\cite{2308.03281} применил многоэтапное контрастивное обучение для построения универсальных текстовых эмбеддингов, включая этапы ненаблюдаемого предобучения и наблюдаемой дообучки. BGE~\cite{2309.07597} предложил комплект ресурсов для развития китайских эмбеддингов общего назначения, включая предобученные модели, корпуса для дообучки и интегрированные кросс-энкодеры.

NV-Embed~\cite{2405.17428} представил архитектурные и методологические инновации для обучения LLM в качестве универсальных моделей эмбеддингов. Ключевые нововведения включают слой латентного внимания для пулинга, удаление каузальной маски внимания во время контрастивного обучения и двухэтапную процедуру обучения с инструкциями. Модели NV-Embed-v1 и NV-Embed-v2 заняли первое место в рейтинге MTEB.

AnglE~\cite{2309.12871} предложил оптимизацию текстовых эмбеддингов в угловом пространстве, что позволяет эффективнее разделять семантически близкие и далёкие тексты. Gecko~\cite{2403.20327} применил дистилляцию знаний из больших языковых моделей для создания компактных, но выразительных текстовых эмбеддингов.

ULLME~\cite{2408.03402} разработал унифицированный фреймворк для получения эмбеддингов из LLM с обучением, усиленным генерацией, что позволяет адаптировать произвольную языковую модель для задач построения представлений. Nomic Embed~\cite{2402.01613} представил воспроизводимый эмбеддер для длинных контекстов с открытым кодом.

Matryoshka Representation Learning~\cite{2205.13147} предложил метод обучения, при котором одна модель способна генерировать эмбеддинги различной размерности, допуская адаптивный выбор длины вектора в зависимости от вычислительных ограничений без существенной потери качества.

Среди наиболее современных подходов следует отметить Qwen3 Embedding~\cite{2506.05176}, основанный на фундаментальных моделях серии Qwen3 и продемонстрировавший передовые результаты как в задачах построения эмбеддингов, так и в переранжировании. Llama-Embed-Nemotron-8B~\cite{2511.07025} представляет универсальную модель текстовых эмбеддингов для многоязычных и кросс-языковых задач. KaLM-Embedding-V2~\cite{2506.20923} применил улучшенные методы обучения и курирования данных. Gemini Embedding~\cite{2503.07891} предложил обобщаемые эмбеддинги на основе моделей семейства Gemini. EmbeddingGemma~\cite{2509.20354} сфокусирован на создании мощных, но компактных текстовых представлений. Diffusion-Pretrained Dense and Contextual Embeddings~\cite{2602.11151} исследовал возможности диффузионного предобучения для получения плотных и контекстных эмбеддингов.

В работе GigaEmbeddings~\cite{2510.22369} представлен фреймворк для обучения высокопроизводительных русскоязычных текстовых эмбеддингов на основе иерархического инструктивного дообучения декодерной LLM (GigaChat-3B). Трёхэтапный конвейер, включающий масштабное контрастивное предобучение, дообучку с жёсткими негативными примерами и мультизадачную генерализацию, позволил достичь результатов уровня state-of-the-art на бенчмарке ruMTEB (69.1 среднего балла).

\subsection{Мультимодальные эмбеддинги и поиск}\label{subsec:multimodal_embeddings}

С развитием мультимодальных моделей возникла потребность в эмбеддингах, способных работать одновременно с текстовыми и визуальными данными.

MM-Embed~\cite{2411.02571} предложил подход к универсальному мультимодальному поиску на основе мультимодальных LLM, позволяющий обрабатывать запросы и документы, содержащие произвольные комбинации текста и изображений. MetaEmbed~\cite{2509.18095} представил метод масштабирования мультимодального поиска на этапе инференса с использованием гибкого позднего взаимодействия, что обеспечивает адаптивность к различным типам запросов.

CoMa~\cite{2511.08480} предложил эффективную парадигму предобучения мультимодальных эмбеддингов <<сжатие, затем сопоставление>>, что позволяет существенно снизить вычислительные затраты. MCA~\cite{2510.15543} сфокусирован на осведомлённости о составе модальностей для робастного композитного мультимодального поиска. MRMR~\cite{2510.09510} представил реалистичный мультидисциплинарный бенчмарк экспертного уровня для оценки мультимодального поиска, требующего глубокого рассуждения.

\subsection{Обзоры и бенчмарки эмбеддингов}\label{subsec:embedding_surveys}

Систематизация результатов в области текстовых эмбеддингов осуществлена в ряде обзорных работ. MTEB (Massive Text Embedding Benchmark)~\cite{2210.07316} предложил стандартизованный бенчмарк из 56 задач для всесторонней оценки моделей текстовых эмбеддингов, ставший де-факто стандартом в сообществе. Обзор \cite{2412.09165} представил всестороннее исследование на стыке текстовых эмбеддингов и больших языковых моделей, систематизировав методы обучения, архитектурные решения и области применения. Работа \cite{2406.01607} содержит анализ наиболее успешных методов на бенчмарке MTEB, выявляя ключевые факторы качества. Сравнительное исследование диффузионных и авторегрессионных языковых моделей с точки зрения текстовых эмбеддингов проведено в~\cite{2505.15045}.

%% =================================================================
\section{Переранжирование}\label{sec:reranking}
%% =================================================================

Переранжирование (reranking) является критически важным этапом конвейера информационного поиска, обеспечивающим уточнение результатов первичного этапа извлечения. Развитие данного направления прошло путь от эвристических методов до современных подходов на основе больших языковых моделей.

\subsection{Обзоры и основы переранжирования}\label{subsec:reranking_surveys}

Комплексный обзор эволюции моделей переранжирования представлен в~\cite{2512.16236}, охватывающий развитие от эвристических методов до подходов на основе LLM. Эмпирический анализ современных моделей переранжирования на основе LLM проведён в~\cite{2508.16757}, где исследуются сильные и слабые стороны различных подходов. Законы масштабирования для переранжирования в информационном поиске изучены в~\cite{2603.04816}, что позволяет прогнозировать качество моделей в зависимости от вычислительных ресурсов.

\subsection{Модели позднего взаимодействия}\label{subsec:late_interaction}

ColBERT~\cite{2004.12832} представил архитектуру позднего взаимодействия, при которой запрос и документ кодируются независимо с помощью BERT, а затем их сходство вычисляется через эффективный механизм мелкозернистого сопоставления. Данный подход позволяет предвычислить представления документов в офлайн-режиме, что обеспечивает ускорение обработки запросов на два порядка по сравнению с классическими BERT-моделями при сохранении конкурентоспособной точности.

ColBERTer~\cite{2203.13088} развил идеи ColBERT, предложив методы контекстуализированного позднего взаимодействия с улучшенным сжатием представлений, что снижает требования к памяти и ускоряет поиск. Jina-ColBERT-v2~\cite{2408.16672} представляет многоязычную модель позднего взаимодействия общего назначения, расширяющую применимость подхода ColBERT на кросс-языковые сценарии.

\subsection{Переранжирование на основе больших языковых моделей}\label{subsec:llm_reranking}

Применение LLM для переранжирования открыло новые возможности благодаря способности моделей к глубокому пониманию релевантности. RankVicuna~\cite{2309.15088} продемонстрировал возможность листового (listwise) переранжирования документов в режиме zero-shot с использованием открытых LLM. RankZephyr~\cite{2312.02724} существенно сократил разрыв в качестве между открытыми и закрытыми моделями, а в ряде случаев превзошёл GPT-4 на задачах переранжирования.

Rank1~\cite{2502.18418} предложил использование дополнительных вычислений на этапе инференса (test-time compute) для улучшения качества переранжирования. Rank-K~\cite{2505.14432} развил эту идею, применив рассуждение на этапе тестирования для листового переранжирования. ListT5~\cite{2402.15838} адаптировал архитектуру Fusion-in-Decoder для листового переранжирования, продемонстрировав улучшение качества zero-shot поиска.

FIRST~\cite{2406.15657} предложил ускоренное листовое переранжирование с использованием декодирования одного токена, что значительно снижает задержку. ProRank~\cite{2506.03487} применил обучение с подкреплением для подготовки промптов малых языковых моделей к переранжированию. RankFlow~\cite{2502.00709} разработал многоролевой коллаборативный рабочий процесс переранжирования с использованием LLM.

Gumbel Reranking~\cite{2502.11116} предложил дифференцируемую оптимизацию переранжирования через Gumbel-softmax, позволяющую обучение модели end-to-end. ReFIT~\cite{2305.11744} использует обратную связь от переранжировщика в процессе инференса для итеративного улучшения результатов поиска. Подход Beyond Sequential Reranking~\cite{2509.07163} продемонстрировал, что направляемый переранжировщиком поиск улучшает результаты на задачах, требующих глубокого рассуждения. Rank-Nexus~\cite{2601.20623} объединил визуальную и текстовую модальности в листовом переранжировании.

%% =================================================================
\section{Генерация с извлечением информации (RAG)}\label{sec:rag}

\begin{figure}[htbp]
    \centering
    \resizebox{0.9\textwidth}{!}{
        \usetikzlibrary{shapes.geometric, arrows.meta, positioning, fit, shadows, shapes.symbols}

\begin{tikzpicture}[
    node distance=1.5cm,
    process/.style={
        draw,
        rectangle,
        rounded corners,
        minimum width=2.5cm,
        minimum height=1cm,
        align=center,
        fill=blue!10,
        drop shadow
    },
    database/.style={
        draw,
        cylinder,
        shape border rotate=90,
        aspect=0.25,
        minimum width=1.5cm,
        minimum height=1.5cm,
        align=center,
        fill=yellow!20,
        drop shadow
    },
    document/.style={
        draw,
        tape,
        tape bend top=none,
        minimum width=2cm,
        minimum height=1cm,
        align=center,
        fill=white,
        drop shadow
    },
    arrow/.style={
        ->,
        >=Stealth,
        thick
    }
]

% Nodes
\node[process, fill=green!10] (query) {User Query};
\node[process, right=of query] (retriever) {Retriever\\(Dense/Sparse)};
\node[database, below=of retriever] (index) {Document\\Index};
\node[document, right=of retriever] (docs) {Top-$k$\\Documents};
\node[process, right=of docs] (augment) {Context\\Augmentation};
\node[process, fill=orange!20, right=of augment] (llm) {LLM};
\node[process, fill=green!10, below=of llm] (answer) {Answer};

% Connections
\draw[arrow] (query) -- (retriever);
\draw[arrow] (retriever) -- (docs);
\draw[arrow] (index) -- (retriever);
\draw[arrow] (docs) -- (augment);
\draw[arrow] (augment) -- (llm);
\draw[arrow] (llm) -- (answer);

% Query also feeds into augmentation — route above intermediate nodes
\draw[arrow] (query.north) -- ++(0,0.8) -| (augment.north);

% Labels
\node[below=0.2cm of index, font=\small] {Knowledge Base};

\end{tikzpicture}

    }
    \caption{Схема работы генерации с извлечением информации (RAG).}
    \label{fig:rag_pipeline}
\end{figure}
%% =================================================================

Генерация с извлечением информации (Retrieval-Augmented Generation, RAG) --- парадигма, объединяющая извлечение релевантных документов из внешних источников с генеративными возможностями языковых моделей. RAG позволяет преодолеть ограничения LLM, связанные с устареванием знаний, галлюцинациями и недостатком доменной экспертизы.

\subsection{Обзоры и основы RAG}\label{subsec:rag_surveys}

Основополагающая работа Lewis et al.~\cite{2005.11401} предложила архитектуру RAG для решения знаниеёмких задач NLP, объединив предобученную seq2seq модель (параметрическая память) с плотным векторным индексом Википедии (непараметрическая память). Модели RAG установили новые результаты state-of-the-art на нескольких бенчмарках для ответов на вопросы в открытом домене.

Комплексный обзор Gao et al.~\cite{2312.10997} систематизировал развитие парадигм RAG, включая наивный RAG, продвинутый RAG и модульный RAG, а также детально рассмотрел компоненты извлечения, генерации и аугментации. Обзор графовых подходов к RAG~\cite{2501.13958} представил систематический анализ GraphRAG --- новой парадигмы, использующей графовое представление знаний для преодоления ограничений традиционного RAG. Обзор контекстного сжатия в RAG~\cite{2409.13385} исследовал методы эффективной компрессии извлечённых документов. Обзор мультимодального RAG~\cite{2504.08748} систематизировал подходы к извлечению и генерации с использованием данных различных модальностей.

\subsection{Основные методы RAG}\label{subsec:core_rag}

Dense Passage Retrieval (DPR)~\cite{2004.04906} показал, что плотные представления, полученные простым двухэнкодерным фреймворком, могут значительно превзойти классические разреженные методы (TF-IDF, BM25) на задачах поиска пассажей, улучшив точность top-20 на 9--19\% абсолютных.

REALM~\cite{2002.08909} впервые продемонстрировал возможность ненаблюдаемого предобучения латентного извлекателя знаний совместно с языковой моделью, используя маскированное языковое моделирование как обучающий сигнал и обратное распространение через этап извлечения, рассматривающий миллионы документов.

Fusion-in-Decoder (FiD)~\cite{2007.01282} предложил архитектуру, в которой несколько извлечённых пассажей независимо кодируются энкодером, а затем объединяются в декодере, что позволяет эффективно агрегировать информацию из множества источников.

Self-RAG~\cite{2310.11511} представил фреймворк саморефлексивной генерации с извлечением, обучающий единую модель адаптивно извлекать пассажи по мере необходимости и генерировать специальные токены рефлексии, что обеспечивает контролируемость поведения модели на этапе инференса. Модели Self-RAG (7B и 13B) значительно превзошли ChatGPT и Llama2-chat с извлечением на ряде задач.

RAPTOR~\cite{2401.18059} предложил рекурсивное абстрактивное процессирование для древовидного извлечения: текстовые фрагменты рекурсивно встраиваются, кластеризуются и суммаризируются, формируя дерево с различными уровнями абстракции. Бенчмарк RGB~\cite{2309.01431} предложил систематическую оценку LLM в контексте RAG.

\subsection{Продвинутые и агентные подходы}\label{subsec:advanced_rag}

Агентный RAG (Agentic RAG) представляет следующий эволюционный этап, встраивающий автономные AI-агенты в конвейер RAG. Обзор~\cite{2501.09136} систематизировал данное направление, описав агентные паттерны проектирования: рефлексию, планирование, использование инструментов и мультиагентное сотрудничество. Систематизация знаний (SoK) по агентному RAG~\cite{2603.07379} предложила таксономию архитектур, методы оценки и направления исследований.

RAG-Gym~\cite{2502.13957} представил среду для систематической оптимизации языковых агентов в контексте RAG. FAIR-RAG~\cite{2510.22344} предложил верный адаптивный итеративный метод уточнения для RAG, обеспечивающий высокое качество генерации. ChunkRAG~\cite{2410.19572} разработал метод фильтрации чанков на уровне LLM, позволяющий отсеивать нерелевантный контекст. Know3-RAG~\cite{2505.12662} представил знание-ориентированный фреймворк RAG с адаптивным извлечением, генерацией и фильтрацией.

\subsection{Графовые подходы к RAG}\label{subsec:graph_rag}

Интеграция графов знаний с RAG позволяет учитывать структурированные связи между сущностями. KG-R1~\cite{2509.26383} предложил эффективный и переносимый агентный подход к RAG на графах знаний, обучаемый с помощью обучения с подкреплением. KG-Infused RAG~\cite{2506.09542} исследовал аугментацию корпусного RAG внешними графами знаний для улучшения качества ответов. TOBUGraph~\cite{2412.05447} разработал подход к извлечению на основе графов знаний, превосходящий стандартный RAG.

\subsection{Безопасность и робастность RAG}\label{subsec:rag_security}

Включение внешних и непроверенных данных в контекст LLM создаёт дополнительные векторы атак. SafeRAG~\cite{2501.18636} предложил бенчмарк безопасности RAG, классифицирующий атаки на четыре типа: <<серебряный шум>>, межконтекстные конфликты, <<мягкая реклама>> и отказ в обслуживании (white DoS). Эксперименты на 14 компонентах RAG продемонстрировали значительную уязвимость существующих систем ко всем типам атак.

Работа~\cite{2402.07179} исследовала влияние пертурбаций промптов на качество генерации в RAG-системах. Riddle Me This~\cite{2502.00306} представил скрытные атаки вывода о принадлежности (membership inference) для RAG, демонстрирующие угрозы приватности в системах извлечения.

%% =================================================================
\section{Системы обеспечения безопасности ИИ}\label{sec:ai_safety}
%% =================================================================

Развёртывание больших языковых и мультимодальных моделей в промышленных приложениях требует надёжных систем модерации контента и обеспечения безопасности. В данном разделе рассмотрены основные подходы к созданию таких систем.

Llama Guard~\cite{2312.06674} --- одна из первых систем безопасности на основе LLM, предназначенная для классификации входных промптов и выходных ответов в контексте диалогов <<человек--ИИ>>. Модель, построенная на Llama2-7b и дообученная на инструкциях, обеспечивает многоклассовую классификацию рисков с возможностью настройки таксономии. Llama Guard 3 Vision~\cite{2411.10414} расширил данный подход на мультимодальный сценарий, обеспечивая безопасность в диалогах, включающих анализ изображений. LlamaFirewall~\cite{2505.03574} представляет систему с открытым кодом для построения безопасных AI-агентов, интегрирующую множество защитных механизмов.

ShieldGemma~\cite{2407.21772} --- комплексный набор моделей модерации контента на основе Gemma2, обеспечивающий предсказание рисков безопасности по ключевым типам вреда (сексуально откровенный контент, опасный контент, домогательства, разжигание ненависти) как во входных данных, так и в генерируемых ответах. ShieldGemma продемонстрировал превосходство над Llama Guard (+10.8\% AU-PRC на публичных бенчмарках). ShieldGemma~2~\cite{2504.01081} расширил данный подход, обеспечив робастную модерацию изображений.

Aegis2.0~\cite{2501.09004} представил разнообразный датасет и таксономию рисков для выравнивания защитных ограждений LLM. NeMo Guardrails~\cite{2310.10501} предложил инструментарий для построения контролируемых и безопасных приложений LLM с программируемыми ограждениями, позволяющий разработчикам задавать правила поведения модели на специализированном языке.

garak~\cite{2406.11036} --- фреймворк для тестирования безопасности LLM, автоматизирующий поиск уязвимостей через систематические атаки-зонды. OmniGuard~\cite{2512.02306} предложил унифицированные омни-модальные ограждения с механизмом обдуманного рассуждения. Qwen3Guard~\cite{2510.14276} представляет модель безопасности на основе архитектуры Qwen3, обеспечивающую классификацию рисков для различных сценариев использования.

%% =================================================================
\section{Уязвимости больших языковых и мультимодальных моделей}\label{sec:vulnerabilities}

\begin{figure}[htbp]
    \centering
    \resizebox{0.9\textwidth}{!}{
        \usetikzlibrary{trees, shadows, positioning, shapes.geometric, arrows.meta}

\begin{tikzpicture}[
    level 1/.style={sibling distance=8cm, level distance=2cm},
    level 2/.style={sibling distance=3.2cm, level distance=2cm},
    every node/.style={
        draw,
        rectangle,
        rounded corners,
        align=center,
        top color=white,
        bottom color=blue!10,
        drop shadow,
        font=\small,
        minimum height=0.8cm,
        inner sep=5pt
    },
    edge from parent/.style={
        draw,
        ->,
        thick,
        >=Stealth
    },
    edge from parent path={
        (\tikzparentnode.south) -- ++(0,-0.5cm) -| (\tikzchildnode.north)
    }
]

\node[bottom color=blue!20, font=\small\bfseries] {Attacks on\\Multimodal LLMs}
    child {node {Jailbreaks}
        child {node {Text-based}}
        child {node {Image-based}}
        child {node {Cross-modal}}
    }
    child {node {Adversarial\\Attacks}
        child {node {Perturbations}}
        child {node {Universal\\Patches}}
    }
    child {node {Prompt\\Injection}
        child {node {Direct}}
        child {node {Indirect}}
    };

\end{tikzpicture}

    }
    \caption{Классификация атак на мультимодальные большие языковые модели.}
    \label{fig:attack_taxonomy}
\end{figure}
%% =================================================================

Несмотря на впечатляющие возможности современных LLM и мультимодальных LLM (MLLM), они остаются уязвимыми к различным типам атак, направленных на обход встроенных механизмов безопасности. Данный раздел систематизирует основные направления исследований в области уязвимостей и методов их эксплуатации.

\subsection{Обзоры уязвимостей}\label{subsec:vulnerability_surveys}

JailbreakZoo~\cite{2407.01599} предоставил обширный обзор методов взлома (jailbreaking) LLM и визуально-языковых моделей (VLM), систематизировав атаки в семь категорий и описав соответствующие защитные стратегии. Обзор <<From LLMs to MLLMs>>~\cite{2406.14859} исследовал ландшафт мультимодального взлома, выявив ограничения текущих подходов и перспективные направления для повышения робастности MLLM.

Unbridled Icarus~\cite{2404.05264} представил обзор потенциальных угроз, связанных с визуальными входами в MLLM, систематизировав риски, привносимые изображениями. Обзор состязательной робастности MLLM~\cite{2503.13962} охватил методы атак и защит в контексте мультимодальных архитектур. Работа~\cite{2307.16680} исследовала ландшафт доверительности генеративных моделей, включая вопросы надёжности, справедливости и безопасности. Комплексный обзор~\cite{2601.03594} рассмотрел механизмы, методы оценки и унифицированные защиты от взлома LLM и VLM.

\subsection{Атаки взлома на визуально-языковые модели}\label{subsec:jailbreak_attacks}

Атаки взлома (jailbreak) направлены на обход механизмов безопасности моделей путём создания специально сконструированных входных данных, вынуждающих модель генерировать запрещённый контент. В мультимодальном контексте визуальная модальность открывает дополнительные поверхности атаки.

CAMO (Cross-Modal Adversarial Multimodal Obfuscation)~\cite{2506.16760} предложил фреймворк атаки типа <<чёрный ящик>>, декомпозирующий вредоносные промпты на семантически безопасные визуальные и текстовые фрагменты, которые восстанавливают вредоносную инструкцию через кросс-модальное рассуждение модели. UltraBreak~\cite{2602.01025} разработал подход к универсальным и переносимым атакам взлома на VLM, стремясь к максимальной обобщаемости атак между моделями.

ImgTrojan~\cite{2403.02910} продемонстрировал возможность взлома VLM с использованием единственного изображения, внедряя вредоносную инструкцию в визуальное представление. MML Attack~\cite{2412.00473} предложил атаку через мультимодальную связь: вредоносная информация распределяется между модальностями, а модель вынуждена интегрировать её для генерации опасного ответа.

BAP (Bi-Modal Adversarial Prompt)~\cite{2406.04031} одновременно оптимизирует состязательные возмущения в обеих модальностях для повышения эффективности атаки. JailBound~\cite{2505.19610} исследовал взлом внутренних границ безопасности VLM через манипуляции в латентном пространстве. SI-Attack~\cite{2501.04931} предложил атаку через нарушение согласованности при перемешивании входных данных.

IDEATOR~\cite{2411.00827} продемонстрировал, что VLM могут быть использованы для взлома самих себя, генерируя собственные вредоносные промпты. AutoJailbreak~\cite{2407.16686} применил LLM для автоматической генерации атак на GPT-4V. PBI-Attack~\cite{2412.05892} предложил управляемую приорами бимодальную интерактивную атаку для максимизации токсичности.

Benign-to-Toxic~\cite{2505.21556} исследовал индуцирование вредоносных ответов из безобидных промптов. Sequential Comics~\cite{2510.15068} использовал последовательные комиксы как средство структурированного визуального повествования для взлома MLLM. Jailbreaks via Multimodal Reasoning~\cite{2601.22398} исследовал атаки через мультимодальное рассуждение.

JaiLIP~\cite{2509.21401} предложил взлом VLM через управляемое потерями возмущение изображений. Defense2Attack~\cite{2509.12724} продемонстрировал, что слабые защиты могут быть использованы для усиления атак. Исследование~\cite{2407.15211} систематически изучило трудности переноса состязательных изображений между различными VLM. The Great Contradiction Showdown~\cite{2410.01438} проанализировал противоречие между скрытностью и эффективностью атак взлома.

\subsection{Состязательные атаки на мультимодальные модели}\label{subsec:adversarial_attacks}

Состязательные атаки (adversarial attacks) направлены на создание входных данных, вызывающих ошибочное поведение модели, при этом возмущения часто остаются незаметными для человека.

Работа~\cite{2502.07987} представила универсальную состязательную атаку на выровненные (aligned) MLLM, генерируя единый возмущающий шаблон, эффективный против множества моделей. AnyDoor~\cite{2402.08577} предложил атаки бэкдоров на этапе тестирования MLLM, не требующие модификации параметров модели. FOA-Attack~\cite{2505.21494} разработал метод атаки на закрытые MLLM через оптимальное выравнивание признаков.

Типографические атаки представляют отдельный класс угроз: работа~\cite{2402.00626} показала, что VLM способны <<обманывать>> сами себя с помощью самогенерируемых типографических атак, а~\cite{2502.08193} расширил данное исследование на сценарий с множественными изображениями.

Состязательная робастность для визуального граундинга MLLM исследована в~\cite{2405.09981}. Атаки на робототехнические VLA (Vision-Language-Action) модели рассмотрены в~\cite{2506.03350}. Переносимость состязательных атак из изображений в видео в контексте MLLM изучена в~\cite{2501.01042}. Visual Stitching~\cite{2506.03614} продемонстрировал, что VLM способны агрегировать разрозненные обучающие патчи, что создаёт дополнительный вектор атаки.

\subsection{Механизмы защиты мультимодальных моделей}\label{subsec:defense_mechanisms}

Разработка эффективных защитных механизмов является одной из наиболее актуальных задач в области безопасности MLLM. Существующие подходы можно классифицировать по нескольким направлениям: преобразование входных данных, робастификация энкодеров, детекция атак и усиление выравнивания.

ECSO (Eyes Closed, Safety On)~\cite{2403.09572} предложил обучение-свободный подход к защите MLLM, адаптивно преобразующий небезопасные изображения в текстовые описания для активации встроенных механизмов безопасности LLM. Данный метод обеспечил улучшение безопасности на 37.6\% на бенчмарке MM-SafetyBench.

BlueSuffix~\cite{2410.20971} разработал метод защиты <<синей команды>>, включающий визуальный очиститель, текстовый очиститель и генератор суффиксов, обучаемый с помощью обучения с подкреплением для повышения кросс-модальной робастности. Sim-CLIP+~\cite{2409.07353} предложил робастный энкодер для защиты VLM от взломов и состязательных атак. Q-MLLM~\cite{2511.16229} применил векторное квантование для обеспечения робастной безопасности MLLM.

SmoothGuard~\cite{2510.26830} использует шумовые возмущения и кластерную агрегацию для защиты MLLM. PSA-VLM~\cite{2411.11543} предложил прогрессивное выравнивание через концептуальные <<бутылочные горлышки>> (concept bottlenecks) для повышения безопасности VLM.

В области детекции атак работа~\cite{learning-to-detect-unknown-jailbreak-attacks-in-large-vision-language-models} исследовала обнаружение неизвестных атак взлома на VLM. Rethinking Jailbreak Detection~\cite{rethinking-jailbreak-detection-of-large-vision-language-models-with-representati} предложил контрастивную оценку на уровне представлений для детекции взломов. E$^2$AT~\cite{multimodal-jailbreak-defense-via-dynamic-joint-optimization-for-mllms} представил мультимодальную защиту через динамическую совместную оптимизацию.

The Safety Reminder~\cite{2506.15734} предложил мягкий промпт для реактивации отложенной осведомлённости о безопасности в VLM. ShiftDC~\cite{2502.13095} исследовал искажение восприятия безопасности в VLM и методы его коррекции. DefenSee~\cite{2512.01185} представил многоракурсный защитный конвейер для мультимодальных взломов. Internal Activation Revision~\cite{2501.16378} предложил защиту VLM без обновления параметров через ревизию внутренних активаций.

Механистическое исследование работы защит от взлома и их ансамблирования проведено в~\cite{2502.14486}. SIA~\cite{2507.16856} предложил повышение безопасности VLM через осведомлённость о намерениях. BPO (Blind Preference Optimization)~\cite{2503.14189} представил метод создания безвредных мультимодальных ассистентов через оптимизацию предпочтений без визуальных данных. Dyslexify~\cite{2508.20570} разработал механистическую защиту от типографических атак в CLIP.

\subsection{Красная команда для мультимодальных моделей}\label{subsec:red_teaming}

Красная команда (red teaming) --- систематическая методология тестирования моделей на уязвимости путём имитации действий злоумышленника.

RTVLM~\cite{2401.12915} представил методологию красной команды для визуально-языковых моделей, систематизировав сценарии тестирования. Arondight~\cite{2407.15050} предложил стандартизированный фреймворк красной команды для VLM с автоматической генерацией мультимодальных промптов взлома: визуальные промпты генерируются VLM из <<красной команды>>, а текстовые --- LLM под управлением агента обучения с подкреплением. Arondight достиг среднего показателя успешности атак 84.5\% на GPT-4.

OpenRT~\cite{2601.01592} представил фреймворк красной команды с открытым кодом для MLLM. Исследование Red Teaming GPT-4V~\cite{2404.03411} проверило устойчивость GPT-4V к одномодальным и мультимодальным атакам взлома. Работа~\cite{2509.15478} исследовала вред, причиняемый через различные модальности промптов. MUSE~\cite{2603.02482} представил платформу для унифицированной оценки безопасности мультимодальных LLM. VERA-V~\cite{2510.17759} предложил вариационный фреймворк инференса для взлома VLM.

\subsection{Атаки внедрения промптов}\label{subsec:prompt_injection}

Атаки внедрения промптов (prompt injection) направлены на манипуляцию поведением LLM-интегрированных приложений путём внедрения вредоносных инструкций в контекст модели.

HouYi~\cite{2306.05499} предложил метод атаки внедрения промптов на LLM-интегрированные приложения по аналогии с традиционными веб-инъекциями, включающий три компонента: предконструированный промпт, инъекцию, вызывающую разделение контекста, и вредоносную полезную нагрузку. Тестирование на 36 реальных приложениях выявило 31 уязвимое.

Работа~\cite{2403.04957} представила единый фреймворк для понимания целей атак внедрения промптов и автоматический градиентный метод генерации универсальных данных для атак. Ранняя категоризация атак внедрения промптов представлена в~\cite{an-early-categorization-of-prompt-injection-attacks-on-large-language-models}.

BIPIA~\cite{2312.14197} предложил бенчмарк для оценки и защиты от непрямых атак внедрения промптов. StruQ~\cite{2402.06363} разработал защиту через структурированные запросы, разделяющие пользовательские инструкции от данных. SecAlign~\cite{2410.05451} применил оптимизацию предпочтений для защиты от внедрения промптов.

WebInject~\cite{2505.11717} исследовал атаки внедрения промптов на веб-агентов. Работа~\cite{2511.01634} оценила устойчивость LLM к внедрению промптов как возникающей угрозе. UniGuardian~\cite{2502.13141} предложил унифицированную защиту для обнаружения внедрения промптов, атак бэкдоров и состязательных атак.

\subsection{Выравнивание безопасности мультимодальных моделей}\label{subsec:safety_alignment}

Выравнивание безопасности (safety alignment) направлено на обучение моделей отказываться от генерации вредоносного контента при сохранении полезности.

Работа~\cite{2405.13581} представила комплексный подход к выравниванию безопасности VLM. MM-RLHF~\cite{2502.10391} предложил следующий шаг в мультимодальном выравнивании через обучение с подкреплением на основе обратной связи от человека. DREAM~\cite{2504.18053} разработал метод декомпозиции рисков для повышения безопасности MLLM через наблюдаемое дообучение и итеративное RLAIF, продемонстрировав улучшение на 16.17\% по метрике safe\&effective относительно GPT-4V.

SEA~\cite{2502.12562} предложил ресурсо-эффективное выравнивание безопасности MLLM через синтетические эмбеддинги. Equilibrate RLHF~\cite{equilibrate-rlhf-towards-balancing-helpfulness-safety-trade-off-in-large-languag} исследовал баланс между полезностью и безопасностью в LLM.

XSafety~\cite{2310.00905} исследовал многоязычную безопасность LLM, продемонстрировав, что модели проявляют неравномерную безопасность на различных языках, что создаёт дополнительные риски. Гипотеза поверхностного выравнивания безопасности~\cite{2410.10862} показала, что существующие механизмы безопасности могут быть более хрупкими, чем представляется, действуя лишь на поверхностном уровне.

\subsection{Бенчмарки и оценка безопасности}\label{subsec:safety_benchmarks}

Систематическая оценка безопасности MLLM требует комплексных бенчмарков, покрывающих различные типы рисков и сценарии атак.

MM-SafetyBench~\cite{mm-safetybench-a-benchmark-for-safety-evaluation-of-multimodal-large-language-mo} --- один из первых бенчмарков для оценки безопасности MLLM, охватывающий основные категории вреда. SafeBench~\cite{2410.18927} предложил комплексный фреймворк для оценки безопасности MLLM, включающий автоматическую генерацию датасета из 2300 вредоносных мультимодальных пар запросов и протокол оценки на основе <<жюри присяжных>> --- коллаборативных LLM.

MultiTrust~\cite{2406.07057} представил всесторонний бенчмарк доверительности MLLM, оценивающий модели по множеству измерений. MOSSBench~\cite{2406.17806} исследовал проблему чрезмерной чувствительности (oversensitivity) мультимодальных моделей к безопасным запросам. VLSBench~\cite{2411.19939} выявил проблему <<визуальной утечки>> в мультимодальных бенчмарках безопасности, когда модели могут определить вредоносность запроса исключительно по текстовой компоненте.

How Many Unicorns~\cite{how-many-unicorns-are-in-this-image-a-safety-evaluation-benchmark-for-vision-llm} представил бенчмарк оценки безопасности VLM. MLLMGuard~\cite{2406.07594} предложил многомерный набор для оценки безопасности MLLM. SaLAD~\cite{2601.04043} исследовал сценарии, в которых MLLM-помощники становятся источниками опасности в повседневной жизни.

Lingua-SafetyBench~\cite{2601.22737} --- бенчмарк для оценки безопасности многоязычных VLM, исследующий влияние языка на проявление уязвимостей. OmniSafeBench-MM~\cite{2512.06589} представил унифицированный бенчмарк и инструментарий для оценки мультимодальных атак взлома и защит. VLDBench~\cite{2502.11361} оценивает мультимодальную дезинформацию с учётом регуляторных требований. PII-VisBench~\cite{2601.05739} фокусируется на безопасности персонально идентифицируемой информации в VLM. T2VSafetyBench~\cite{2407.05965} оценивает безопасность генеративных моделей <<текст-видео>>.

\subsection{Безопасность аудио- и видеомодальностей}\label{subsec:audio_video_safety}

Помимо визуальной модальности, аудио- и видеовходы создают дополнительные поверхности атаки на мультимодальные модели.

Работа~\cite{2410.23861} исследовала уязвимости аудио-мультимодальных LLM, продемонстрировав, что аудиомодальность является <<ахиллесовой пятой>> данных систем. SpeechGuard~\cite{speechguard-exploring-the-adversarial-robustness-of-multimodal-large-language-mo} исследовал состязательную робастность мультимодальных LLM в контексте речевых входов. Работа~\cite{breaking-audio-large-language-models-by-attacking-only-the-encoder} показала возможность взлома аудио-LLM через атаку исключительно на энкодер.

Переносимость состязательных атак в видео-MLLM через кросс-модальный подход <<изображение-видео>> исследована в~\cite{2501.01042}, демонстрируя, что состязательные примеры, созданные для изображений, могут быть эффективно перенесены на видеомодели. Атаки вывода о принадлежности (membership inference) на VLM рассмотрены в~\cite{membership-inference-attacks-against-large-vision-language-models}.

%% =================================================================

В рамках данного обзора литературы систематизированы основные направления исследований, связанные с тематикой настоящей диссертации. Анализ показывает стремительное развитие как методов атак на мультимодальные модели, так и защитных механизмов, при этом разрыв между атакующими и защитными возможностями остаётся значительным. Результаты автора диссертации в области русскоязычных языковых моделей (GigaChat~\cite{2506.09440}) и текстовых эмбеддингов (GigaEmbeddings~\cite{2510.22369}) непосредственно связаны с рассмотренными направлениями и обеспечивают основу для дальнейших исследований безопасности в мультимодальном контексте.
